\section{Dijkstra}
\whenuse{Caminos m\'inimos desde una fuente con pesos no negativos. Para pesos
0/1 usar 0--1 BFS; con negativos, Bellman--Ford.}

Distancias desde una fuente con pesos no negativos.

\begin{lstlisting}[style=cpp]
vector<ll> dijkstra(const vector<vector<pair<int,ll>>> &graph, int source) {
    int n = graph.size();
    vector<ll> distance(n, LINF);
    priority_queue<pair<ll,int>, vector<pair<ll,int>>,
                   greater<pair<ll,int>>> pending;
    distance[source] = 0;
    pending.push({0, source});

    while (!pending.empty()) {
        auto [dist_v, v] = pending.top(); pending.pop();
        if (dist_v != distance[v]) continue;
        for (auto [to, weight] : graph[v]) {
            if (weight < 0) throw invalid_argument("negative edge");
            if (distance[to] > dist_v + weight) {
                distance[to] = dist_v + weight;
                pending.push({distance[to], to});
            }
        }
    }
    return distance;
}
vector<WeightedEdge> edges = {
    {0, 1, 4},
    {0, 2, 5},
    {1, 2, -2},
    {2, 3, 3}
};
auto [distance, negative_cycle] = bellman_ford(4, edges, 0);
if (negative_cycle) cout << "Hay ciclo negativo\n";
else cout << distance[3] << '\n'; // 5: 0 -> 1 -> 2 -> 3
\end{lstlisting}

\complexity{$O((V+E)\log V)$; memoria $O(V+E)$.}

\section{0--1 BFS}
\whenuse{Caminos m\'inimos cuando cada arista pesa exactamente 0 o 1; evita el
factor logar\'itmico de Dijkstra usando una deque.}

Para pesos exclusivamente cero o uno.

\begin{lstlisting}[style=cpp]
vector<int> zero_one_bfs(const vector<vector<pair<int,int>>> &graph,
                         int source) {
    int n = graph.size();
    vector<int> distance(n, INF);
    deque<int> q;
    distance[source] = 0; q.push_front(source);
    while (!q.empty()) {
        int v = q.front(); q.pop_front();
        for (auto [to, weight] : graph[v]) {
            assert(weight == 0 || weight == 1);
            if (distance[to] > distance[v] + weight) {
                distance[to] = distance[v] + weight;
                if (weight == 0) q.push_front(to);
                else q.push_back(to);
            }
        }
    }
    return distance;
}
vector<vector<pair<int,int>>> graph(4);
graph[0].push_back({1, 0}); // 0 -> 1, costo 0
graph[0].push_back({2, 1}); // 0 -> 2, costo 1
graph[1].push_back({2, 0}); // 1 -> 2, costo 0
graph[2].push_back({3, 1}); // 2 -> 3, costo 1

vector<int> distance = zero_one_bfs(graph, 0);
cout << distance[3] << '\n'; // 1
\end{lstlisting}

\complexity{$O(V+E)$ tiempo y memoria.}

\section{Bellman--Ford}
\whenuse{Caminos m\'inimos con pesos negativos y detecci\'on de ciclos negativos
alcanzables; s\'olo para grafos moderados por su costo cuadr\'atico.}

Admite pesos negativos. Tras $V-1$ rondas, una relajaci\'on adicional indica
un ciclo negativo alcanzable desde la fuente.

\begin{lstlisting}[style=cpp]
struct WeightedEdge { int from, to; ll weight; };

pair<vector<ll>, bool> bellman_ford(
        int n, const vector<WeightedEdge> &edges, int source) {
    vector<ll> distance(n, LINF);
    distance[source] = 0;
    for (int iteration = 0; iteration < n - 1; ++iteration) {
        bool changed = false;
        for (auto [from, to, weight] : edges) {
            if (distance[from] == LINF) continue;
            i128 raw = (i128)distance[from] + weight;
            ll candidate = raw < -LINF ? -LINF
                         : raw > LINF ? LINF : (ll)raw;
            if (candidate < distance[to]) {
                distance[to] = candidate;
                changed = true;
            }
        }
        if (!changed) break;
    }
    bool negative_cycle = false;
    for (auto [from, to, weight] : edges)
        if (distance[from] != LINF &&
            (i128)distance[from] + weight < distance[to])
            negative_cycle = true;
    return {distance, negative_cycle};
}
\end{lstlisting}

\complexity{$O(VE)$ tiempo y $O(V)$ memoria.}

\section{Floyd--Warshall}
\whenuse{Distancias entre todos los pares cuando $V$ es peque\~no; tambi\'en para
clausura transitiva o detectar ciclos negativos mediante la diagonal.}

Todos contra todos. Inicializar diagonal en cero, aristas en su peso m\'inimo y
el resto en \texttt{LINF}.

\begin{lstlisting}[style=cpp]
void floyd_warshall(vector<vector<ll>> &distance) {
    int n = distance.size();
    for (int k = 0; k < n; ++k)
        for (int i = 0; i < n; ++i) if (distance[i][k] != LINF)
            for (int j = 0; j < n; ++j) if (distance[k][j] != LINF)
                distance[i][j] = min(distance[i][j],
                    distance[i][k] + distance[k][j]);
}

bool negative_cycle_between(const vector<vector<ll>> &d, int a, int b) {
    for (int k = 0; k < (int)d.size(); ++k)
        if (d[a][k] != LINF && d[k][k] < 0 && d[k][b] != LINF)
            return true;
    return false;
}
\end{lstlisting}

\complexity{$O(V^3)$ tiempo y $O(V^2)$ memoria.}
